% Partie : théorie de la démonstration
% Chapitre : recherche de preuves
% Section : complétude

\documentclass[../../../include/open-logic-section]{subfiles}

\begin{document}

\olfileid[fr]{pt}{ps}{cpl}

\olsection{Complétude}

Nous montrons maintenant que, si l'algorithme s'arrête
sans succès ou ne s'arrête jamais, le séquent initial
$\Gamma \Sequent \Delta$ n'a pas de preuve.
Nous construisons pour cela une structure~$\Struct M$
telle que $\Sat{M}{!A}$ pour tout $!A \in \Gamma$,
et $\Sat/{M}{!A}$ pour tout~$!A \in \Delta$.
Toutes les formules de~$\Gamma$ y sont donc vraies
et toutes celles de~$\Delta$ fausses : le séquent
$\Gamma \Sequent \Delta$ est invalide. Par correction
de \Log{G3c}, $\Gamma \Sequent \Delta$ n'a alors
aucune preuve. Nous conservons les conventions
du chapitre : formules closes et absence d'égalité logique.

La structure $\Struct{M}$ est définie à partir de
deux ensembles de formules $\Theta, \Xi$ et d'un
ensemble~$C$ de constantes. On obtient $\Theta, \Xi$
et~$C$ sur une \emph{branche d'échec} de la recherche.
Il s'agit d'une suite de séquents
$\Pi_n \Sequent \Lambda_n$, munis d'ensembles de
constantes~$C_n$, qui commence par le séquent final
$\Pi_0 \Sequent \Lambda_0$, égal à $\Gamma \Sequent \Delta$.
Chaque séquent $\Pi_{n+1} \Sequent \Lambda_{n+1}$
est une prémisse de $\Pi_n \Sequent \Lambda_n$
engendrée à l'étape~$n+1$, et aucun séquent
de la branche n'est une feuille de clôture.
On note $C_n$ l'ensemble de constantes associé
à~$\Pi_n \Sequent \Lambda_n$.

Si l'algorithme s'arrête sans succès sur un séquent
supérieur ne contenant que des formules atomiques,
le chemin menant à ce séquent est une branche
d'échec finie. Sinon, lorsque la recherche ne
s'arrête jamais, elle engendre des arbres de
plus en plus grands. Leur union est un arbre
infini à ramifications finies. Chaque nœud
a au plus deux prémisses immédiates.
Le lemme de K\H{o}nig fournit une branche infinie.
Comme aucune feuille de clôture n'est développée,
cette branche ne passe par aucune clôture ;
c'est une branche d'échec.

À partir de $\Pi_n \Sequent \Lambda_n$ et $C_n$,
on pose $\Theta = \bigcup_n \Pi_n$ et
$\Xi = \bigcup_n \Lambda_n$, en supprimant les
indices et les répétitions de formules, ainsi que
$C = \bigcup_n C_n$. Pour une branche finie,
ces unions portent sur les seuls séquents
effectivement présents.

On définit alors~$\Struct{M}$ comme suit.
\begin{enumerate}
  \item Son domaine $\Domain{M}$ est l'ensemble
  des termes clos construits à partir de~$C$
  et des symboles de fonction qui figurent dans
  $\Gamma \Sequent \Delta$, s'il y en a.
  Ce domaine est non vide, puisque l'ensemble
  initial de constantes est non vide.

  \item Pour chaque constante $c \in C$,
  on pose $\Assign{c}{M} = c$.

  \item Si $f$ est un symbole de fonction d'arité $m$
  présent dans $\Gamma \Sequent \Delta$, et
  $t_1$, \dots, $t_m \in \Domain{M}$, on pose
  $\Assign{f}{M}(t_1, \dots, t_m) =
  f(t_1, \dots, t_m)$.

  \item Si $R$ est un symbole de prédicat d'arité $m$
  présent dans $\Gamma \Sequent \Delta$, on pose
  $\Assign{R}{M} = \Setabs{\tuple{t_1, \dots, t_m} \in
  \Domain{M}^m}{R(t_1, \dots, t_m) \in \Theta}$.
\end{enumerate}

La structure $\Struct{M}$ est un \emph{modèle de termes} :
son domaine est constitué de termes, et l'interprétation
des constantes et des fonctions garantit que
$\Value{t}{M} = t$. Les formules atomiques de~$\Theta$
définissent $\Assign{R}{M}$ de façon que
$\Sat{M}{R(t_1, \dots, t_m)}$ équivaille à
$R(t_1, \dots, t_m) \in \Theta$.
Les symboles du langage qui ne figurent pas dans
le séquent initial, hormis les constantes déjà
interprétées, peuvent recevoir des interprétations
arbitraires sur ce domaine non vide. Cela ne
modifie la valeur d'aucune formule de la branche.

\begin{lem}\ollabel{lem:truth}
Pour tout $!A \in \Theta \cup \Xi$ :
\begin{enumerate}
  \item si $!A \in \Theta$, alors $\Sat{M}{!A}$ ;
  \item si $!A \in \Xi$, alors $\Sat/{M}{!A}$.
\end{enumerate}
\end{lem}

\begin{proof}
On raisonne par récurrence sur $\depth{!A}$.

Supposons d'abord $!A$ atomique : soit
$!A \ident \lfalse$, soit $!A \ident R(t_1, \dots, t_m)$.
Aucun antécédent $\Pi_n$ de la branche d'échec
$\Pi_n \Sequent \Lambda_n$ ne contient $\lfalse$,
car un tel séquent serait un axiome.
D'autre part, $\Sat{M}{R(t_1, \dots, t_m)}$ équivaut
à $R(t_1, \dots, t_m) \in \Theta$ par définition
de~$\Struct{M}$. Ainsi, si $!A \in \Theta$,
alors $\Sat{M}{!A}$.

Les occurrences atomiques persistent sur la branche :
si $!A \in \Pi_n$, alors $!A \in \Pi_{n+1}$,
et si $!A \in \Lambda_n$, alors $!A \in \Lambda_{n+1}$,
tant qu'un séquent suivant existe.
Si $!A \in \Theta \cap \Xi$, on aurait donc,
à un même stade~$n$, $!A \in \Pi_n \cap \Lambda_n$.
Le séquent $\Pi_n \Sequent \Lambda_n$ serait
une feuille de clôture, ce qui est exclu.
On a donc $!A \notin \Theta \cap \Xi$
pour une formule atomique~$!A$. Si $!A \in \Xi$, alors
$!A \notin \Theta$, d'où $\Sat/{M}{!A}$
par l'interprétation des prédicats.
Le cas de l'absurde, qui est toujours faux,
satisfait lui aussi cette conclusion.

Pour l'étape de récurrence, supposons les deux
assertions établies pour toutes les formules
de profondeur inférieure à celle de~$!A$.
Nous distinguons les opérateurs principaux.

L'équité garantit qu'une occurrence non atomique
$!A^i$ présente dans $\Pi_n \Sequent \Lambda_n$
est finalement réduite. Elle persiste dans
$\Pi_{n+1} \Sequent \Lambda_{n+1}$ tant qu'elle
n'est pas choisie. À chaque réduction, l'indice
minimal \emph{parmi les occurrences non atomiques}
augmente, puisque les indices introduits sont
strictement plus grands que tous les anciens.
Les atomes peuvent garder de petits indices
sans intervenir dans ce choix. On atteint donc
un séquent $\Pi_n \Sequent \Lambda_n$ où $!A^i$
est l'occurrence non atomique d'indice~$i$ minimal,
puis on la réduit à l'étape $n+1$.
Ainsi, si $!A \in \Theta$, il existe un stade $n$
où $\Pi_n = \Pi_n', !A^i$ et où cette occurrence
est choisie ; si $!A \in \Xi$, il existe de même
un stade avec $\Lambda_n = \Lambda_n', !A^i$.

\begin{enumerate}
  \item Supposons $!A \in \Theta$ et
  $!A \ident !B \land !C$. Pour un certain $n$,
  $\Pi_n = \Pi_n', (!B \land !C)^i$.
  Le séquent $\Pi_{n+1} \Sequent \Lambda_{n+1}$
  est la prémisse de \LeftR{\land}, avec
  $\Pi_{n+1} = \Pi_n', !B^k, !C^{k+1}$.
  Donc $!B, !C \in \Theta$.
  Par hypothèse de récurrence, $\Sat{M}{!B}$
  et $\Sat{M}{!C}$. Par définition de
  $\Sat{M}{}$, on obtient $\Sat{M}{!B \land !C}$.

  \item Supposons $!A \in \Xi$ et
  $!A \ident !B \land !C$. Pour un certain $n$,
  $\Lambda_n = \Lambda_n', !B \land !C$.
  Le séquent $\Pi_{n+1} \Sequent \Lambda_{n+1}$
  est l'une des deux prémisses de \RightR{\land}
  dont la conclusion est
  $\Pi_n \Sequent \Lambda'_n, !B \land !C$.
  Ainsi, $\Lambda_{n+1} = \Lambda_n', !B^k$
  ou $\Lambda_{n+1} = \Lambda_n', !C^k$.
  Donc $!B \in \Xi$ ou $!C \in \Xi$.
  Par hypothèse de récurrence, $\Sat/{M}{!B}$
  ou $\Sat/{M}{!C}$ ; par définition de
  $\Sat{M}{}$, il vient $\Sat/{M}{!B \land !C}$.

  \item Les cas $!A \ident \lnot !B$,
  $!A \ident !B \lor !C$ et $!A \ident !B \lif !C$
  sont analogues. Leurs détails sont laissés
  en exercice. Les conditions de vérité
  à vérifier sont données ci-dessous.

  \item Supposons $!A \in \Theta$ et
  $!A \ident \lexists[x][!B(x)]$.
  Pour un certain $n$,
  $\Pi_n = \Pi_n', \lexists[x][!B(x)]^i$.
  Le séquent $\Pi_{n+1} \Sequent \Lambda_{n+1}$
  est la prémisse de \LeftR{\lexists} :
  $\Pi_{n+1} = \Pi_n', !B(c)^k$
  pour une certaine constante~$c$.
  Donc $!B(c) \in \Theta$ et~$c \in C$.
  Par hypothèse de récurrence, $\Sat{M}{!B(c)}$.
  Par définition de $\Sat{M}{}$,
  on obtient $\Sat{M}{\lexists[x][!B(x)]}$.

  \item Supposons $!A \in \Xi$ et
  $!A \ident \lexists[x][!B(x)]$.
  Pour un certain $n$,
  $\Lambda_n = \Lambda_n', \lexists[x][!B(x)]^i$.
  Chaque réduction de $\lexists[x][!B(x)]$
  à droite conserve $\lexists[x][!B(x)]$
  dans la prémisse, avec un nouvel indice.
  Une branche d'échec contenant
  $\lexists[x][!B(x)]$ à droite est donc infinie :
  elle ne peut finir sur une feuille atomique.
  Par équité, cette occurrence est réduite
  une infinité de fois.

  Pour tout terme $t \in \Domain{M}$, ses
  constantes, en nombre fini, appartiennent
  toutes à~$C_n$ à partir d'un certain stade.
  Le terme a un rang fini dans l'énumération ;
  les seuls termes susceptibles de le précéder
  sont donc en nombre fini. L'historique
  conservé lors des réindexations assure que
  ce terme finit par être choisi pour réduire
  $\lexists[x][!B(x)]$ à droite.
  Pour tout $t \in \Domain{M}$, il existe ainsi
  un stade~$n$ où $!B(t) \in \Lambda_n$ ;
  par conséquent, $!B(t) \in \Xi$.
  Par hypothèse de récurrence,
  $\Sat/{M}{!B(t)}$ pour tout $t \in \Domain{M}$.
  Par définition de $\Sat{M}{}$,
  on conclut $\Sat/{M}{\lexists[x][!B(x)]}$.

  \item Les cas $!A \ident \lforall[x][!B(x)]$
  sont analogues ; leurs détails sont laissés
  en exercice.
\end{enumerate}

Les cas restants se lisent dans les règles.
Sur la branche, une négation
$\lnot !B$ dans $\Theta$ place $!B$ dans $\Xi$,
et une négation dans $\Xi$ place $!B$ dans $\Theta$.
Une disjonction $!B \lor !C$ dans $\Theta$
place au moins l'une de ses composantes dans
$\Theta$, tandis qu'une disjonction dans $\Xi$
place les deux dans $\Xi$.
Une implication $!B \lif !C$ dans $\Theta$
place $!B$ dans $\Xi$ ou $!C$ dans $\Theta$ ;
dans $\Xi$, elle place $!B$ dans $\Theta$
et $!C$ dans $\Xi$.
L'hypothèse de récurrence donne chaque fois
la condition de vérité ou de fausseté voulue.

Une universelle $\lforall[x][!B(x)]$ dans
$\Theta$ est conservée et réduite une infinité
de fois. Le même argument d'énumération donne
$!B(t) \in \Theta$ pour tout terme du domaine ;
elle est donc vraie. Une universelle dans
$\Xi$ fournit une instance $!B(c) \in \Xi$
avec $c \in C$ ; elle est donc fausse.
Les instances ont une profondeur strictement
inférieure à celle de la formule quantifiée,
ce qui justifie l'emploi de l'hypothèse de
récurrence.
%
\end{proof}

\begin{cor}\ollabel{cor:complete}
L'algorithme de recherche pour \Log{G3c} est complet :
s'il s'arrête sans succès ou ne s'arrête jamais,
le séquent initial~$\Gamma \Sequent \Delta$
est invalide et n'a donc pas de preuve.
\end{cor}

\begin{proof}
Une telle recherche possède une branche d'échec,
à partir de laquelle on définit $\Struct{M}$.
Par le \olref{lem:truth}, pour tout $!A \in \Gamma$,
on a $\Sat{M}{!A}$, et pour tout $!A \in \Delta$,
on a $\Sat/{M}{!A}$. En effet, $\Pi_0 = \Gamma$
et $\Lambda_0 = \Delta$, donc
$\Gamma \subseteq \Theta$ et $\Delta \subseteq \Xi$.
Le séquent $\Gamma \Sequent \Delta$ est invalide.
Par correction de~\Log{G3c},
$\Gamma \Sequent \Delta$ n'a aucune preuve.
\end{proof}

\begin{cor}
Le calcul \Log{G3c} est complet :
si $\Gamma \Sequent \Delta$ est valide,
alors $\Log{G3c} \Proves \Gamma \Sequent \Delta$.
\end{cor}

\begin{proof}
On raisonne par contraposition. Supposons
$\Log{G3c} \Proves/ \Gamma \Sequent \Delta$.
L'algorithme ne peut s'arrêter avec succès,
puisqu'aucune preuve n'existe ; il s'arrête
donc sans succès ou poursuit indéfiniment
la recherche. D'après le \olref{cor:complete},
$\Gamma \Sequent \Delta$ est invalide.
\end{proof}

\end{document}
